\section*{第 \ref{ch:b1:plate-kinematics} 章——球面上的板块运动学}

\begin{solution}{exo:b1:plate-kinematics:1}
(a) $\qty{0.210}{\degree\per\Ma} = 0.210\times\pi/180 = \qty{3.665e-3}{rad}$
每百万年，即 \qty{3.665e-9}{rad.\text{年}^{-1}}；取
$\qty{1}{\text{年}} = \qty{3.156e7}{s}$，则为 \qty{1.161e-16}{rad.s^{-1}}。(b) $360/0.210 =
\qty{1714}{\Ma}$，超过地球年龄的三分之一。(c) $\omega R =
\num{3.665e-9}\times\qty{6.371e9}{mm} = \qty{23.35}{mm.\text{年}^{-1}}$。
\end{solution}

\begin{solution}{exo:b1:plate-kinematics:2}
$v = \omega R\sin\Delta = 23.35\sin\Delta$：在 \ang{30} 处为 \qty{11.7}{mm.\text{年}^{-1}}，
在 \ang{60} 处为 \qty{20.2}{mm.\text{年}^{-1}}，在 \ang{90} 处为 \qty{23.35}{mm.\text{年}^{-1}}。扩张
速度在边界离极点最远处最快，直至
\ang{90}：位于欧亚板块与北美板块之间，在边界的南端，
靠近亚速尔群岛，距极点 \ang{78}。
\end{solution}

\begin{solution}{exo:b1:plate-kinematics:3}
该亚\omterm{def:b1:plate-kinematics:magnetic-anomaly}{极性时}持续了 $1.072 - 0.988 = \qty{0.084}{\Ma}$，半速率以
\unit{mm.\text{年}^{-1}} 表示，也可写作 \unit{km\per\Ma}。(a) $50\times0.084 =
\qty{4.2}{km}$（速率为 \qty{50}{mm.\text{年}^{-1}}）；$10\times0.084 = \qty{0.84}{km}$（速率为
\qty{10}{mm.\text{年}^{-1}}）。(b) 条带横跨洋脊轴线：
$2\times10\times0.781 = \qty{15.6}{km}$。
\end{solution}

\begin{solution}{exo:b1:plate-kinematics:4}
半速率为 $66/2.581 = \qty{25.6}{mm.\text{年}^{-1}}$（向西）与
$64/2.581 = \qty{24.8}{mm.\text{年}^{-1}}$（向东）；全速率为
$130/2.581 = \qty{50.4}{mm.\text{年}^{-1}}$。扩张
几乎对称：西翼增长快约
\qty{3}{\percent}。
\end{solution}

\begin{solution}{exo:b1:plate-kinematics:5}
$\cos\Delta = \sin\ang{61.8}\sin\ang{85} + \cos\ang{61.8}\cos\ang{85}
\cos(\ang{0} - \ang{139.6}) = 0.8813\times0.9962 + 0.4726\times0.08716\times
(-0.7615) = 0.87795 - 0.03137 = 0.84658$，故 $\Delta = \ang{32.16}$，
$v = 23.35\sin\ang{32.16} = \qty{12.4}{mm.\text{年}^{-1}}$。这略高于
实测的 \qtyrange{11}{12}{mm.\text{年}^{-1}}，但\omterm{def:b1:plate-kinematics:euler-pole}{欧拉矢量}是对
453 个速率的单一最佳拟合，各速率平均偏离 \qty{0.7}{mm.\text{年}^{-1}}，
且慢速的 Gakkel 洋脊附近的个别速率，
属于精度最低的一类。
\end{solution}

\begin{solution}{exo:b1:plate-kinematics:6}
(a) $\cos\Delta = \sin\ang{61.8}\sin\ang{39.44} + \cos\ang{61.8}\cos\ang{39.44}
\cos(\ang{169.43}) = 0.20111$：故 $\Delta = \ang{78.40}$。(b) $v = 23.35
\sin\ang{78.40} = \qty{22.87}{mm.\text{年}^{-1}}$，与交点正北处实测的
\qty{23}{mm.\text{年}^{-1}} 相符。(c) 按照方法中的公式，分子为
$\sin\ang{169.43}\cos\ang{61.8} = 0.08668$，分母为
$\cos\ang{39.44}\sin\ang{61.8} - \sin\ang{39.44}\cos\ang{61.8}\cos\ang{169.43}
= 0.97573$：极点位于 $\beta = \ang{5.08}$，几乎正北，且
欧亚板块向 $\ang{95.1}$ 方向运动。(d) 转换断层与
相对运动平行：其\omterm{def:b1:geological-maps:attitude}{走向}约为 \ang{095}，几乎为东西向。
\end{solution}

\begin{solution}{exo:b1:plate-kinematics:7}
(a) $\vec\omega = \omega(\cos\varphi_p\cos\lambda_p, \cos\varphi_p\sin\lambda_p,
\sin\varphi_p)$：$\vec\omega_{\mathrm{EU\,NA}} = (-0.0756, 0.0643, 0.1851)$，
$\vec\omega_{\mathrm{NB\,NA}} = (0.0333, 0.0282, 0.2289)$，单位为
\unit{\degree\per\Ma}。(b) $\vec\omega_{\mathrm{NB\,EU}} =
\vec\omega_{\mathrm{NB\,NA}} - \vec\omega_{\mathrm{EU\,NA}} = (0.1089, -0.0361,
0.0438)$，长度为
\[
  \sqrt{0.011859 + 0.001303 + 0.001918} = \qty{0.1228}{\degree\per\Ma}.
\]
其极点位于\omterm{def:b1:spherical-astronomy:geographic}{纬度}
$\arcsin(0.0438/0.1228) = \ang{20.9}$\,N，\omterm{def:b1:spherical-astronomy:geographic}{经度}
$\arctan(-0.0361/0.1089) = \ang{-18.3}$，即西经 \ang{18.3}（
$x$ 分量为正，故不加 \ang{180}）。(c) 努比亚板块与欧亚板块之间的边界
直接测量很少：在亚速尔群岛与
直布罗陀之间，仅有 Gloria 断层的\omterm{def:b1:spherical-astronomy:horizontal}{方位角}。横跨大西洋的两个扩张运动
各有数百个速率测量值，它们的差值
继承了这些测量的精度。
\end{solution}

\begin{solution}{exo:b1:plate-kinematics:8}
(a) 绕 $\vec e_z$ 转 \ang{90}：$\cos\ang{90} = 0$，$\vec e_z\cdot\vec e_x = 0$
且 $\vec e_z\times\vec e_x = \vec e_y$，故 $\vec e_x$ 转到 $\vec e_y$。再
绕 $\vec e_x$ 转 \ang{90}：$\vec e_x\cdot\vec e_y = 0$，且
$\vec e_x\times\vec e_y = \vec e_z$，故 $\vec e_y$ 转到 $\vec e_z$：该点
最终到达北极。(b) 先绕 $\vec e_x$ 转：$\vec e_x$ 位于轴上，
$\vec e_x\times\vec e_x = \vec 0$，且 $(1 - \cos\ang{90})(\vec e_x\cdot\vec e_x)
\vec e_x = \vec e_x$，故它不动；再绕 $\vec e_z$ 转，则到达
$\vec e_y$，即北纬 \ang{0}、东经 \ang{90}。(c) 两种结果不同：\omterm{def:b1:plate-kinematics:finite-rotation}{有限
旋转}不满足交换律。连续的重建必须通过
按正确顺序乘以旋转矩阵来复合，而不能通过把角度
或极点相加来完成。
\end{solution}

\begin{solution}{exo:b1:plate-kinematics:9}
(a) 在速度空间中，把 $a$ 置于原点，$b$ 位于
东方 \qty{30}{mm.\text{年}^{-1}} 处，$c$ 位于南方 \qty{40}{mm.\text{年}^{-1}} 处。(b) $\vec v_{CB} =
\vec v_{CA} - \vec v_{BA} = (0, -40) - (30, 0) = (-30, -40)$（东、北分量）：
\qty{50}{mm.\text{年}^{-1}}，指向\omterm{def:b1:spherical-astronomy:horizontal}{方位角}
$\arctan\bigl((-30)/(-40)\bigr) + \ang{180} = \ang{216.9}$，即西南方向。(c)
洋脊垂直于该方向：其\omterm{def:b1:geological-maps:attitude}{走向}为
$\ang{216.9} - \ang{90} = \ang{126.9}$，即西北—东南向。
\end{solution}

\begin{solution}{exo:b1:plate-kinematics:10}
(a) 在 Rodrigues 矩阵中，单位矩阵贡献 $3\cos\alpha$，矩阵
$\vec n\,\vec n^{\mathsf T}$ 的迹为 $n_x^2 + n_y^2 + n_z^2 = 1$，乘以
$1 - \cos\alpha$，反对称矩阵则不贡献：$\operatorname{tr}R =
3\cos\alpha + 1 - \cos\alpha = 1 + 2\cos\alpha$。(b) 此时 $\operatorname{tr}R = 0$，
故 $\cos\alpha = -1/2$，$\alpha = \ang{120}$。转轴满足 $R\vec n =
\vec n$：$R(x, y, z) = (z, x, y)$，故 $x = y = z$，且 $\vec n =
(1, 1, 1)/\sqrt3$，位于\omterm{def:b1:spherical-astronomy:geographic}{纬度} $\arcsin(1/\sqrt3) = \ang{35.26}$\,N，
东经 \ang{45}。(c) $R\vec e_x = \vec e_y$，$R\vec e_y = \vec e_z$，
$R\vec e_z = \vec e_x$：三个坐标轴被轮换，正如绕
对角线转动三分之一圈所产生的效果。Rodrigues 公式取 $\alpha = +\ang{120}$、转轴为 $\vec n$ 时
确实给出 $\vec e_x\mapsto -\tfrac12\vec e_x + \tfrac12(1, 1, 1) +
\tfrac{\sqrt3}{2}\cdot\tfrac{1}{\sqrt3}(0, 1, -1) = (0, 1, 0)$，故旋转
从北纬 \ang{35.26}、东经 \ang{45} 上方俯视为逆时针方向。
\end{solution}

\begin{solution}{exo:b1:plate-kinematics:11}
(a) 在转换断层上，相对速度 $\vec\omega\times\vec r$ 沿
断层方向，故与 $\vec t$ 平行；它与 $\vec\omega$ 垂直，因此
$\vec t\cdot\vec\omega = 0$。过球心且垂直于 $\vec t$ 的平面
包含 $\vec r$（因 $\vec t$ 水平）和 $\vec\omega$：它是
过该点且垂直于断层的\omterm{def:b1:spherical-astronomy:spherical-triangle}{大圆}所在的平面，
极点位于此\omterm{def:b1:spherical-astronomy:spherical-triangle}{大圆}上。两个这样的平面交于
$\vec t_1\times\vec t_2$ 所在的直线。(b) 在节点处 $\vec t_1 = (0.5436, 0.8367, -0.0673)$，
在冰岛 $\vec t_2 = (0.4942, 0.8639, -0.0968)$，故
$\vec t_1\times\vec t_2 = (-0.02285, 0.01936, 0.05612)$，长度为 $0.06362$。
其\omterm{def:b1:spherical-astronomy:geographic}{纬度}为 $\arcsin(0.05612/0.06362) = \ang{61.9}$\,N，其\omterm{def:b1:spherical-astronomy:geographic}{经度}为
$\ang{180} - \arctan(0.01936/0.02285) = \ang{139.7}$\,E（$x$ 分量为负）：
即本章的极点。(c) 精度很差：两个\omterm{def:b1:spherical-astronomy:spherical-triangle}{大圆}以很小的夹角相交，
因为两者的\omterm{def:b1:geological-maps:attitude}{走向}在边界的 \ang{26} 范围内仅相差 \ang{8}。
\omterm{def:b1:geological-maps:attitude}{走向}改变一度，极点即移动约 \ang{15}。实际
测定会将许多转换断层与扩张速率结合起来，扩张速率沿边界的变化
决定了到极点的角距离。
\end{solution}

\begin{solution}{exo:b1:plate-kinematics:12}
(a) $(94 - 4)/3 = \qty{30}{mm.\text{年}^{-1}}$。(b) 相隔一年的两次读数之差为
$24 - 4 = 20$、$25 - 5 = 20$、$26 - 6 = 20$、$65 - 15 = 50$、$74 - 24 = 50$，
$75 - 25 = 50$、$76 - 26 = 50$、$85 - 65 = 20$ 及 $94 - 74 = 20$（mm）。排序后：
五个 20 与四个 50；中位数为 \qty{20}{mm.\text{年}^{-1}}。(c) 相隔一年的两个
位置处于季节振荡的同一位相，
其差中相互抵消。跳跃只影响跨越它的配对，
此处九对中有四对跨越。只要跨越跳跃的配对不足一半，
中位数便不受其影响，而 (a) 中的端点估计把整个跳跃
当作运动，结果偏高 \qty{50}{\percent}。该方法在跳跃
过于频繁、以致大多数一年期配对都含有跳跃时，或测站
运动不平稳（例如大地震之后）时失效。
\end{solution}

\begin{solution}{pb:b1:plate-kinematics:1}
\textbf{1.} $\omega = 0.210\times(\pi/180)/\num{e6} = \qty{3.665e-9}{rad.\text{年}^{-1}}$；
$\omega R = \num{3.665e-9}\times\qty{6.371e9}{mm} = \qty{23.35}{mm.\text{年}^{-1}}$。

\textbf{2.} $\cos\Delta = \sin\ang{61.8}\sin\ang{64.5} + \cos\ang{61.8}
\cos\ang{64.5}\cos(\ang{-157.6}) = 0.79546 - 0.18809 = 0.60737$，故
$\Delta = \ang{52.60}$。

\textbf{3.} $v = 23.35\sin\ang{52.60} = 23.35\times0.7944 =
\qty{18.55}{mm.\text{年}^{-1}}$。

\textbf{4.} 分子为 $\sin\ang{157.6}\cos\ang{61.8} = 0.18011$，而
分母为
\[
  \cos\ang{64.5}\sin\ang{61.8} - \sin\ang{64.5}\cos\ang{61.8}\cos\ang{157.6} =
  0.37940 + 0.39436 = 0.77376;
\]
两者均为正，故 $\beta = \arctan 0.2328 =
\ang{13.1}$。欧亚板块向 $\ang{13.1} + \ang{90} = \ang{103.1}$ 方向运动。

\textbf{5.} 向东：$18.55\sin\ang{103.1} = \qty{18.07}{mm.\text{年}^{-1}}$；向北：
$18.55\cos\ang{103.1} = \qty{-4.20}{mm.\text{年}^{-1}}$，即
向南 \qty{4.20}{mm.\text{年}^{-1}}。

\textbf{6.} $\cos\Delta = 0.8851$，$\Delta = \ang{27.74}$，$v = 23.35
\sin\ang{27.74} = \qty{10.9}{mm.\text{年}^{-1}}$。两板块绕一个极点转动，该极点
位于西伯利亚东部，Gakkel 海岭靠近它，速率随
$\sin\Delta$ 减小，尽管两板块相同。

\textbf{7.} 向西：$94/9.779 = \qty{9.61}{mm.\text{年}^{-1}}$；向东：$97/9.779 =
\qty{9.92}{mm.\text{年}^{-1}}$。

\textbf{8.} $(94 + 97)/9.779 = 191/9.779 = \qty{19.53}{mm.\text{年}^{-1}}$。

\textbf{9.} 由 C2An：$54.1/2.581 = \qty{20.96}{mm.\text{年}^{-1}}$；由
C1n 底部：$19.8/0.781 = \qty{25.35}{mm.\text{年}^{-1}}$。两者并不一致：越年轻的
边界，速率越高。真实的速率变化是可能的，但
距离上存在一个在各年龄上都相同的\omterm{def:b1:plate-kinematics:triple-junction}{稳定}偏差，恰好会得到同样的结果。

\textbf{10.} 向西与向东的距离之比为 0.96 比 0.98：东翼
增长多出几个百分点。\omterm{def:b1:plate-kinematics:spreading-rate}{全扩张速率}仅取决于两距离
之和，故不对称性并不改变它。

\textbf{11.} $(191 - 5.0)/9.779 = \qty{19.02}{mm.\text{年}^{-1}}$；
$(54.1 - 5.0)/2.581 = \qty{19.02}{mm.\text{年}^{-1}}$；$(19.8 - 5.0)/0.781 =
\qty{18.95}{mm.\text{年}^{-1}}$。三者现在彼此相差
不超过 \qty{0.1}{mm.\text{年}^{-1}}。同样的 \qty{5}{km} 占最年轻一对
\qty{19.8}{km} 的四分之一，却不到最古老一对 \qty{191}{km} 的
\qty{3}{\percent}；古老边界给出的速率更可靠。

\textbf{12.} 修正后的\omterm{def:b1:plate-kinematics:spreading-rate}{半扩张速率}为 $186/9.779/2 = \qty{9.51}{mm.\text{年}^{-1}}$，
而 $9.51\times0.084 = \qty{0.80}{km}$，约为
每条边界所偏移的 \qty{2.5}{km} 的三分之一。人们会预期如此窄的条带会被模糊
掉。它在 Reykjanes 海岭上可以识别，而在
缺乏岩浆的 Gakkel 海岭上则不存在，这表明充足的岩浆供应
即使在窄条带中也能忠实地记录磁场。

\textbf{13.} 修正后的速率为 18.95 至 \qty{19.02}{mm.\text{年}^{-1}}，与
板块矢量给出的 \qty{19.0}{mm.\text{年}^{-1}} 相符；若不作修正，则
高出 3\% 至 33\%。

\textbf{14.} $\cos\Delta = \sin\ang{64.139}\sin\ang{59.018} +
\cos\ang{64.139}\cos\ang{59.018}\cos\ang{27.554} = 0.97054$，故 $\Delta =
\ang{13.94}$；弦长为 $2R\sin(\Delta/2) = 12742\times\sin\ang{6.97} =
\qty{1546}{km}$。

\textbf{15.} $1000/17.5 = 57$ 年。

\textbf{16.} 两方向相差 $\ang{123.1} - \ang{114.8} =
\ang{8.3}$：沿\omterm{def:b1:spherical-astronomy:spherical-triangle}{大圆}的水平分量为
$18.9\cos\ang{8.3} = \qty{18.70}{mm.\text{年}^{-1}}$，沿弦的分量为
$18.70\cos\ang{6.97} = \qty{18.56}{mm.\text{年}^{-1}}$，约为 \qty{18.6}{mm.\text{年}^{-1}}。

\textbf{17.} 弦长每年增加 \qty{17.5}{mm.\text{年}^{-1}}，
比预测的 \qty{18.6}{mm.\text{年}^{-1}} 少 \qty{6}{\percent}。两个测站都远离
板块边界，因此这一差异并非局地性的。要么板块今天的运动
比过去三百万年的平均速度略慢，
要么两种估计之一带有较小的系统误差；仅凭这些测量
无法判断。仅由卫星定位计算得到的\omterm{def:b1:plate-kinematics:euler-pole}{欧拉矢量}
对该弦的预测值为 \qty{17.1}{mm.\text{年}^{-1}}，与实测值接近。

\textbf{18.} 东向：$21.12 - (-1.07) = \qty{22.19}{mm.\text{年}^{-1}}$；北向：
$-6.97 - (-0.97) = \qty{-6.00}{mm.\text{年}^{-1}}$。速率 $\sqrt{22.19^2 + 6.00^2} =
\qty{23.0}{mm.\text{年}^{-1}}$，方向朝 $\ang{180} - \arctan(22.19/6.00) =
\ang{105.1}$。

\textbf{19.} 方向接近板块矢量的 \ang{103}，但
速率比其 \qty{18.6}{mm.\text{年}^{-1}} 高出四分之一。两个测站都位于边界变形的区域内，
霍芬靠近最大的冰盖，随着冰层变薄，\omterm{def:b1:layered-earth:layers}{地壳}回弹，
因而向上运动
\qty{12.9}{mm.\text{年}^{-1}}，并沿水平方向背离冰盖运动。
位于边界带内的测站测得的是板块运动加上局地变形。

\textbf{20.} $0.210\times10 = \ang{2.10}$。

\textbf{21.} 该点沿半径为 $R\sin\Delta$ 的\omterm{def:b1:spherical-astronomy:spherical-triangle}{小圆}运动：
$6371\times0.7944\times(2.10\times\pi/180) = \qty{185.5}{km}$。经修正后，
自 C5n 年轻边缘以来形成的\omterm{def:b1:layered-earth:layers}{地壳}宽度为 \qty{186}{km}，
历时 \qty{9.8}{Ma}，与之几乎相同：
运动在一千万年间一直保持\omterm{def:b1:plate-kinematics:triple-junction}{稳定}。

\textbf{22.} 分量（东向，北向）：欧亚板块
$(22.9\sin\ang{95.1}, 22.9\cos\ang{95.1}) = (22.81, -2.04)$，努比亚板块
$(19.2\sin\ang{103.8}, 19.2\cos\ang{103.8}) = (18.65, -4.58)$。欧亚板块相对于努比亚板块：
$(4.16, 2.54)$，即 \qty{4.9}{mm.\text{年}^{-1}}，方向
朝 $\arctan(4.16/2.54) = \ang{58.6}$，东北偏东。

\textbf{23.} 三条洋脊的交汇点无论速度如何都是\omterm{def:b1:plate-kinematics:triple-junction}{稳定}的。第 22 题中的速度使
欧亚板块与努比亚板块沿东北偏东至西南偏西方向相互分离：
它们在交汇点处的边界必须张开，且速度缓慢；
这就是特塞拉裂谷，其东南偏东的\omterm{def:b1:geological-maps:attitude}{走向}使张开方向呈斜交。

\textbf{24.} 第一部分：冰岛中部为 \qty{18.6}{mm.\text{年}^{-1}}
（第二部分剖面处为 \qty{19.0}{mm.\text{年}^{-1}}）。第二部分：经修正后为 18.95 至
\qty{19.02}{mm.\text{年}^{-1}}。第三部分：雷克雅未克–斯塔万格弦长
增加 \qty{17.5}{mm.\text{年}^{-1}}，而板块矢量预测为
\qty{18.6}{mm.\text{年}^{-1}}；冰岛境内的测站对给出 \qty{23.0}{mm.\text{年}^{-1}}，
并非板块运动。板块矢量与条带结果一致；
卫星结果相差 \qty{6}{\percent}。条带平均时间跨度可达一千万年，
板块矢量平均时间跨度为 \qty{3.16}{Ma}；卫星测得的
则是最近二十至三十年的运动。

\textbf{25.} 板块矢量给出冰岛境内 \qty{18.55}{mm.\text{年}^{-1}}，
其旁经修正的条带给出 \qty{19.0}{mm.\text{年}^{-1}}：取到每年最接近的毫米数，
冰岛的\omterm{def:b1:plate-kinematics:spreading-rate}{全扩张速率}为
\textbf{\qty[text-series-to-math=true]{19}{mm.\text{年}^{-1}}}。
\end{solution}
